% ---------------------------------------------------------------------------
% Data section draft -- filament CMB lensing with BOSS CMASS pairs
%
% Written to match the checked-in analysis pipeline and its archived outputs.
%
% Bib keys used here are in paper/styau/refs.bib:
%   Eisenstein2011, Dawson2013, Alam2015, Ross2012, Gorski2005, Riebe2013
%
% "% CHECK" marks numbers or statements that still require verification.
% ---------------------------------------------------------------------------

\section{Data}
\label{sec:data}

The measurement combines three ingredients: spectroscopic galaxy positions
from BOSS, which define the pairs and the single-galaxy control; the
corresponding BOSS random catalogs, which trace the angular selection and
measure footprint-dependent offsets; and the Planck CMB lensing convergence
map, which supplies the signal. An $N$-body mock of the CMASS sample and its
convergence field provides a $\Lambda$CDM reference signal
(Section~\ref{sec:sim}).

\subsection{BOSS CMASS galaxies}
\label{sec:data_cmass}

Galaxy positions come from the CMASS sample of the Baryon Oscillation
Spectroscopic Survey \citep[BOSS;][]{Dawson2013}, part of SDSS-III
\citep{Eisenstein2011}, in its twelfth data release
\citep{Alam2015, Reid2016}. CMASS targets were selected with color and
magnitude cuts designed to follow a passively evolving stellar-population
locus, so that the sample has an approximately constant stellar mass across
its redshift range; hence the name. The galaxies are luminous red galaxies
hosted by haloes of typical mass $M\simeq10^{13}\,h^{-1}\mathrm{M}_{\odot}$,
with a satellite fraction of about 10 percent \citep{White2011}. Massive
hosts are useful here because their one-halo lensing signal can be measured
and modeled, and simulations predict that pairs of such haloes are often
connected by filaments \citep{Colberg2005}.

We use both Galactic caps, referred to as North and South, and restrict the
sample to $0.4<z<0.7$. This is slightly wider than the $0.43<z<0.7$ range
adopted for the BOSS clustering analyses \citep{Reid2016}.
% CHECK: document the reason for using 0.4 rather than 0.43.
After this cut the sample contains $792{,}294$ galaxies, $579{,}089$ in the
North and $213{,}205$ in the South, with a summed weight of $847{,}528$.
The effective survey area is approximately $6{,}850\,\mathrm{deg}^{2}$ in the
North and $2{,}525\,\mathrm{deg}^{2}$ in the South, for a total of
$9{,}376\,\mathrm{deg}^{2}$ \citep{Reid2016}. % CHECK against Reid+2016 Table 2
Figure~\ref{fig:footprint_zdist} shows the footprint of both caps together
with the Planck lensing mask. The sample has a median redshift of $z=0.539$
and a weight-averaged redshift of $\langle z\rangle=0.545$; the two caps have
very similar redshift distributions.

Each galaxy carries the standard BOSS systematic weights
\citep{Ross2012, Reid2016}. We combine them as
\begin{equation}
  w = w_{\rm seeing}\, w_{\rm star}\,
      \left(w_{\rm noz} + w_{\rm cp} - 1\right),
  \label{eq:weights}
\end{equation}
where $w_{\rm seeing}$ and $w_{\rm star}$ correct the angular selection
function for variations in seeing and stellar density, $w_{\rm cp}$ accounts
for close pairs lost to fiber collisions, and $w_{\rm noz}$ accounts for
redshift failures. We do not apply the FKP weight because it is designed to
optimize clustering measurements rather than this stacking estimator.
In the single-galaxy stacks we multiply $w$ by a lensing-efficiency weight
$\Sigma_{\rm crit}^{-2}$, where
$\Sigma_{\rm crit}\propto D_{\rm A}(z_{\rm s})/[D_{\rm A}(z)\,D_{\rm A}(z,z_{\rm s})]$
is evaluated with the CMB as the source plane, $z_{\rm s}=1090$, and
Planck 2018 angular-diameter distances. This favors the redshifts at which a
given mass produces the largest convergence. The same weight is applied to
the random points, so that the galaxy and random stacks are weighted
identically before subtraction. In the pair stacks the weight of each pair is
the product $w_{1}w_{2}$ of its members' systematic weights, without the
lensing-efficiency factor. The single-galaxy and pair estimators therefore
have slightly different effective redshift weightings.
% CHECK: quantify the effect of removing the Sigma_crit weight from the
% single-galaxy stacks (weighting_sensitivity.py).

Right ascension and declination are converted to Galactic longitude and
latitude $(l,b)$ to match the native frame of the Planck maps. Redshifts are
converted to comoving distances with the Planck 2018 cosmology
\citep{Planck2018lensing} and expressed in $\hinvMpc$ with $h=0.6766$. The
North and South caps are stacked as a single joint sample: every galaxy
enters with its own weight, so the relative contribution of the two caps is
set by their galaxy content rather than by a separate average
(Section~\ref{sec:method}).

\begin{figure}
  \centering
  \includegraphics[width=0.48\linewidth]{footprint.pdf}
  \includegraphics[width=0.48\linewidth]{zdist.pdf}
  \caption{Left: footprint of the BOSS CMASS North and South samples in
  Galactic coordinates, overlaid on the Planck 2018 lensing analysis mask.
  Right: weighted redshift distribution of the CMASS galaxies in each cap
  and in total, with the $0.4<z<0.7$ selection shaded and the median
  redshift of the selected sample marked.
  }
  \label{fig:footprint_zdist}
\end{figure}

\subsection{Random catalogs}
\label{sec:data_random}

The BOSS random catalogs \citep{Reid2016} sample the angular selection
function of each cap at roughly fifty times the galaxy density, with
redshifts assigned from the observed galaxy distribution. The randoms
therefore share the galaxies' $n(z)$ but contain no clustering. We apply the
same $0.4<z<0.7$ cut and use the full catalogs, $40.9$ million points over
both caps, for every random stack. Replacing the full catalog by a 10 percent
subsample shifts the $10\,\hinvMpc$ bridge excess by
$5\times10^{-4}$, comparable to the measured signal.
% CHECK: combine_filament_jackknife.py docstring.
Random points carry unit systematic weight and, in the single-galaxy stacks,
the same lensing-efficiency weight as the galaxies.

The randoms serve three purposes. First, a stack of the convergence map
around random points measures the mean convergence of the footprint,
including any residual mean-field or mask-related offset, and is subtracted
from the single-galaxy stack. Second, pairs of random points selected with
the same separation criteria as the galaxy pairs provide the equivalent
correction for the pair stack. Because the number of random pairs is of
order $10^{9}$, the full random catalog is needed to reduce the noise in this
correction. Third, a seeded 10 percent subsample of the randoms traces the
survey footprint on a coarse HEALPix grid to define the jackknife regions
used for error estimation (Section~\ref{sec:method}).

\subsection{Planck CMB lensing map}
\label{sec:data_planck}

We use the Planck 2018 lensing reconstruction \citep{Planck2018lensing},
release \texttt{COM\_Lensing\_4096\_R3.00}. The convergence harmonic
coefficients $\kappa_{LM}$ are taken from the minimum-variance combination
of the temperature and polarization quadratic estimators applied to the
SMICA foreground-cleaned maps. The Planck pipeline subtracts the estimator
mean field, and the coefficients are provided up to
$L_{\rm max}=4096$. As discussed in Section~\ref{sec:intro}, the
reconstruction is noise dominated on the small angular scales relevant to
this work.

Before stacking we multiply $\kappa_{LM}$ by a Gaussian beam of
$8\,\mathrm{arcmin}$ full width at half maximum and synthesize a HEALPix map
\citep{Gorski2005} at $N_{\rm side}=2048$, corresponding to a pixel scale of
$1.7\,\mathrm{arcmin}$. This smoothing suppresses small-scale reconstruction
noise. At $z=0.55$, $8\,\mathrm{arcmin}$ spans
$\simeq3.3\,\hinvMpc$ comoving. We test the dependence on smoothing scale in
Section~\ref{sec:method}. We sample the smoothed map at the sky positions
corresponding to the comoving stack grid by nearest-pixel lookup, without
interpolation.

We apply the Planck lensing analysis mask distributed with the release,
which removes the Galactic plane, point sources, and regions of poor
foreground cleaning and retains 67.1 percent of the sky. Masked pixels
receive zero weight in every stack, so that the stacked value at each grid
position is a weighted mean over unmasked pixels only. The CMASS footprint
lies almost entirely within the unmasked region: 96.5 percent of North and
94.7 percent of South galaxies, 96.0 percent overall, fall on unmasked
pixels. The remainder lie mostly in the point-source and cluster holes of the
mask and at the low-latitude edge of the North cap
(Fig.~\ref{fig:footprint_zdist}).
% CHECK: the fraction whose full 100 h^-1 Mpc cutout is unmasked is not yet computed.

\subsection{$N$-body simulation and mock catalog}
\label{sec:sim}

We calculate the $\Lambda$CDM reference signal from the BigMDPL run of the
MultiDark suite \citep{Klypin2016, Riebe2013}, a dark-matter-only simulation
of a periodic box $2500\,\hinvMpc$ on a side with $3840^{3}$ particles of mass
$2.36\times10^{10}\,h^{-1}\mathrm{M}_{\odot}$ in a Planck cosmology with
$\Omega_{\rm m}=0.307$ and $h=0.678$. We use the snapshot at $a=0.6464$,
$z=0.547$, close to the weighted mean redshift of the CMASS sample. Haloes
are taken from the public Rockstar catalog; we keep host haloes only,
discarding subhaloes, and use the Rockstar virial mass $M_{\rm vir}$. The
HOD is applied to hosts with
$\log_{10}(M_{\rm vir}/h^{-1}\mathrm{M}_{\odot})\geq12.5$.

Mock galaxies are placed with a central-only halo occupation distribution of
the form adopted for CMASS by \citet{White2011},
\begin{equation}
  \langle N_{\rm cen}\rangle(M) = \tfrac{1}{2}\left[1+\mathrm{erf}
  \left(\frac{\ln(M/M_{\rm cut})}{\sqrt{2}\,\sigma}\right)\right],
  \label{eq:hod}
\end{equation}
with $\sigma=0.94$ in $\ln M$. We use
$\log_{10}(M_{\rm cut}/h^{-1}\mathrm{M}_{\odot})=13.246$, rather than the
value $13.04$ adopted by \citet{White2011}. With our halo definition and
catalog mass floor, the latter gives a number density of
$4.21\times10^{-4}\,h^{3}\,\mathrm{Mpc}^{-3}$. We set the threshold to match
the CMASS central-galaxy density,
$2.70\times10^{-4}\,h^{3}\,\mathrm{Mpc}^{-3}$: the observed density of
$3\times10^{-4}\,h^{3}\,\mathrm{Mpc}^{-3}$ multiplied by a central fraction
of 0.9.
% Provenance: analysis/sim/results/galaxies_hod_nmatch.csv.meta.json
% (run 2026-07-20, occupation seed 20260720). 13.246 solves nbar = 0.9 * 3e-4
% over hosts_minmass12p5.bin; make_hod_catalog.py still defaults to the
% published White+11 value 13.04, so the paper run passes --log10-mcut.
Each host receives a galaxy with probability $\langle N_{\rm cen}\rangle$.
The input halo catalog ends at the stated mass floor, where the adopted
occupation is still $0.034$, and therefore truncates the low-mass tail of
the model. The mock contains no satellites, whereas about 10 percent of
CMASS galaxies are satellites \citep{White2011}; we return to this limitation
in Section~\ref{sec:discussion}. Halo peculiar velocities are included in
the redshift-space pair selection.

The second halo sample used for the sensitivity tests consists of hosts
within $\pm0.05$ dex of
$\log_{10}(M_{\rm vir}/h^{-1}\mathrm{M}_{\odot})=13.0$
(Section~\ref{sec:method}).

The simulated convergence map is built from the public 1-in-200 particle
subsample, which contains $2.8\times10^{8}$ particles. Each sampled particle
has an effective mass equal to 200 simulation particles. We project the
particles along one box axis through
the full $2500\,\hinvMpc$ depth onto a periodic grid of $0.1\,\hinvMpc$
pixels. If $N_{\rm p}(\boldsymbol{x})$ is the number of particles in a pixel
of width $\ell$, $N_{\rm tot}$ is the total number of projected particles,
and $L$ is the box width, the mean-subtracted surface density is
$\Sigma(\boldsymbol{x})=m_{\rm eff}[N_{\rm p}(\boldsymbol{x})/\ell^{2}
-N_{\rm tot}/L^{2}]$. We convert it to convergence as
$\kappa=\Sigma/\Sigma_{\rm crit}$ with the comoving critical density for a
lens at $z=0.547$ and the CMB as source,
$\Sigma_{\rm crit}=8.88\times10^{14}\,h\,\mathrm{M}_{\odot}\,{\rm Mpc}^{-2}$.
The map is smoothed with a periodic two-dimensional Gaussian of FWHM
$3.31\,\hinvMpc$, which is $8\,\mathrm{arcmin}$ at the comoving distance
$1424.5\,\hinvMpc$ of the snapshot. Because the box is projected at a single
redshift, the mock reproduces the mean signal of structure correlated with
the pairs but not the full line-of-sight integral or the Planck
reconstruction noise; it is used for the reference signal, not for the
observational covariance. The single-galaxy stack from this map has been
checked against an independent direct particle-count estimator; the two
agree to 3 percent inside $10\,\hinvMpc$.

Mock pairs are selected in redshift space from a random 50 percent subsample
of the catalog, which reduces the pair count while preserving the halo-mass
selection in expectation. The projection axis serves as the line of sight,
and each galaxy's coordinate along it is displaced by
$(1+z)\,v_{\parallel}/H(z)$, where $v_{\parallel}$ is its line-of-sight
peculiar velocity, before the line-of-sight cut is applied. This models the
effect of peculiar velocities on the pair selection. Transverse
separations are computed with periodic boundaries.
Nonphysical pairs, close in projection but separated along the line of
sight in real space by $20$--$120\,\hinvMpc$, are selected from the same
catalog. They are not close in three dimensions and therefore provide a null
test of the control.
